% 分类关系按正文独立组织：三组、八个研究问题、五组三项代表课题。
% 分支只表示归属；没有因果箭头，也不按节点面积、顺序编码数值。
\begin{tikzpicture}[figure picture,foundation picture,
 font=\figurefont\fontsize{8.25}{10.5}\selectfont,
 group/.style={draw=FigureStroke,line width=.7pt,rounded corners=1.5pt,minimum height=8mm,minimum width=52mm},
 topic/.style={font=\FigureNoteFont\fontsize{7.75}{10}\selectfont,anchor=west},
 question/.style={anchor=west,inner sep=0pt}]
 \path[use as bounding box] (0,0) rectangle (169,94);
 \node[group,minimum width=86mm] (root) at (84.5,89) {真实使用中的系统何时可以被依赖？};
 \node[group,fill=FigureBlueFill] (G1) at (26,72) {要求与影响};
 \node[group,fill=MutedGreenFill] (G2) at (84.5,72) {证据与行为};
 \node[group,fill=MutedPurpleFill] (G3) at (143,72) {防护与治理};
 \draw[FigureStroke,line width=.65pt] (root.south)--(84.5,81) (26,81)--(143,81);
 \foreach \g/\x in {G1/26,G2/84.5,G3/143}
  \draw[FigureStroke,line width=.65pt] (\x,81)--(\g.north);
 \foreach \x/\g in {3/G1,61.5/G2}{
  \draw[FigureStroke,line width=.65pt] (\x,68)--(\x,8);
  \foreach \y in {59,34,9}\draw[FigureStroke,line width=.65pt] (\x,\y)--({\x+4},\y);
 }
 \draw[FigureStroke,line width=.65pt] (120,68)--(120,34);
 \draw[FigureStroke,line width=.65pt] (120,59)--(124,59) (120,34)--(124,34);
 \node[question] at (8,59) {要求可操作化};
 \node[question] at (8,34) {未知识别与转交};
 \node[question] at (8,9) {风险、收益与负担分配};
 \node[question] at (66.5,59) {变化、适应与恢复};
 \node[question] at (66.5,34) {解释、审计与补救};
 \node[question] at (66.5,9) {目标、监督与权限};
 \node[question] at (125,59) {隐私、安全与行动保护};
 \node[question] at (125,34) {联合保证与持续验证};
 \node[topic] at (8,53) {危害情景};
 \node[topic] at (8,48.5) {可检验主张};
 \node[topic] at (8,44) {处置阈值};
 \node[topic] at (8,28) {校准与覆盖};
 \node[topic] at (8,23.5) {拒答与转交};
 \node[topic] at (8,19) {主张级核验};
 \node[topic] at (66.5,53) {复合偏移};
 \node[topic] at (66.5,48.5) {变化监测};
 \node[topic] at (66.5,44) {适应与恢复};
 \node[topic] at (66.5,28) {归因与反事实};
 \node[topic] at (66.5,23.5) {机制解释};
 \node[topic] at (66.5,19) {补救有效性};
 \node[topic] at (125,53) {数据影响与撤回};
 \node[topic] at (125,48.5) {自适应攻击};
 \node[topic] at (125,44) {最小权限};
\end{tikzpicture}
